%%%%%%%%%%  scan idx 71  |  folio 58  %%%%%%%%%%
\textbf{Corollaire 2.} \emph{Supposons $X$ connexe de dimension $n$, sans singularités en dimension $n-1$. Si $D$ et $D'$ sont deux diviseurs linéairement équivalents sur $X$, on a $\gamma(D)\equiv\gamma(D')\quad\mathrm{mod.}\ \underline K(X)^2$.}

En vertu du corollaire 1, on peut supposer $X$ non singulière~; le corollaire résulte alors de la formule générale~:
\begin{equation}\tag{2.31}
 \gamma(D)=1-\gamma(L(-D))\qquad\mathrm{mod}\ \underline K(X)^2,
\end{equation}
où $L(-D)$ désigne le fibré vectoriel de rang $1$ défini par le diviseur $-D$. Pour vérifier (2.31), on l'écrit sous la forme~:
\begin{equation}\tag{2.31 bis}
 \gamma(L(-D))=1-\gamma(D)\qquad\mathrm{mod}\ \underline K(X)^2,
\end{equation}
et l'on note que les deux membres sont des homomorphismes du groupe des diviseurs dans le groupe multiplicatif des éléments inversibles de l'anneau $\underline K(X)/\underline K(X)^2$ (noter que $\underline K(X)^1/\underline K(X)^2$ est de carré nul, ce qui se vérifie très facilement sur la formule explicite (2.15), en utilisant le corollaire 1). On est ainsi ramené au cas où $D$ est une hypersurface irréductible. Plus généralement, si $D$ est une hypersurface dont les composantes sont disjointes, on a l'égalité~:
\begin{equation}\tag{2.32}
 \gamma(D)=1-\gamma(L(-D))
\end{equation}
comme il résulte de la suite exacte de faisceaux
\[
0\longrightarrow\cO_X(L(-D))\longrightarrow\cO_X\longrightarrow\cO_D\longrightarrow0.
\]

\textbf{Théorème 2.12.} \emph{Soient $X$ une variété quasi-projective non singulière, $Z$ et $Z'$ deux cycles homogènes de codimension $p$ sur $X$ qui sont linéairement équivalents. On a alors~:}
\begin{equation}\tag{2.33}
 \gamma_X(Z)\equiv\gamma_X(Z')\qquad\mathrm{mod}\ \underline K(X)^{p+1}.
\end{equation}
Par hypothèse, on peut trouver un cycle $Y$ de dimension $n-p+1$ sur 
